\section{案例对比与设计模式}

\subsection{PPO / SAC 在不同任务上的表现}

在多项 benchmark 中的规律如下：

\begin{itemize}
    \item \textbf{低维离散 env}（CartPole、Acrobot）：PPO 足以胜任，SAC 的 entropy 可能导致高 exploration variance。
    \item \textbf{高维连续 env}（Humanoid、Walker2D）：SAC 更稳定，SAC + twin Q networks for target smoothing 避免 Q divergence。
    \item \textbf{分布转移场景}（Sim2Real）：Replay buffer + prioritized sampling + policy bootstrap layer 是关键，SAC 的 off-policy 特性更具优势。
\end{itemize}

\begin{importantbox}{混合方案}
在 Sim2Real or offline RL 场景下，常见做法是先用 SAC 收敛 policy，再用 PPO 重新 fine-tune 以满足 safety constraints（PPO 提供切实的 clipping guardrails）。
\end{importantbox}

\subsection{工程设计模式}

两个常见模式：

\begin{enumerate}
    \item \textbf{Dual Agent Pattern}：PPO policy + SAC policy 并行训练，training job 根据 metrics 自动选择更稳定的版本部署。
    \item \textbf{Safety Gate Pattern}：在 deployment path 前置一个 rule-based gate（supervisory policy），在 gate 判定 KL divergence > threshold 时暂停 agent 并 fallback。
\end{enumerate}

\begin{warningbox}{不要只看 reward}
某些 env reward 设计过于简单（如仅完成任务即高 reward），容易让 agent exploit reward shape。务必结合 episode length、entropy、KL divergence 一起评估。
\end{warningbox}

\subsection{本章小结}

案例对比展示了 PPO 在低维 env 中的稳定性、SAC 在高维 env 与 Sim2Real 的优势，并推荐使用 dual agent 与 safety gate 模式来组合能力与可控性。

